\chapter{Análise sintática LL(1)}

\section{Introdução}

Nos capítulos anteriores, apresentamos duas técnicas para análise
sintática descendente: combinadores de parsing e Parsing Expression
Grammars. Ambas constroem uma derivação iniciando pelo símbolo
inicial da gramática e, a cada passo, escolhem uma regra para
expandir o não-terminal mais à esquerda.

Neste capítulo, estudamos outro algoritmo descendente: o
\textbf{analisador LL(1)}. O nome tem três componentes:

\begin{itemize}
  \item
    O primeiro \textbf{L} indica que a entrada é lida da
    \textbf{esquerda para a direita} (\emph{left-to-right}).
  \item
    O segundo \textbf{L} indica que o analisador produz uma
    \textbf{derivação mais à esquerda} (\emph{leftmost derivation}).
  \item
    O \textbf{1} indica que o analisador usa apenas \textbf{um
    símbolo de lookahead} para decidir qual produção aplicar.
\end{itemize}

A ideia central é substituir a recursão por um \textbf{algoritmo
dirigido por tabela}: uma tabela de parsing \(M[A, a]\) indica, para
cada não-terminal \(A\) e cada terminal \(a\) (incluindo o marcador
de fim de entrada \(\$\)), qual produção deve ser aplicada, ou
sinaliza um erro sintático.

O código discutido neste capítulo está em
\begin{flushleft}
  \verb|Parsing/LL/BuildTable.hs| e\\
  \verb|Parsing/LL/Parser.hs|
\end{flushleft}

\section{Exemplo Informal}

Considere a gramática de expressões aritméticas a seguir (sem
recursão à esquerda, fatorada à esquerda):

\[
  \begin{array}{lcl}
    E  & \to & T\; E' \\
    E' & \to & +\; T\; E' \mid \lambda \\
    T  & \to & F\; T' \\
    T' & \to & *\; F\; T' \mid \lambda \\
    F  & \to & (\; E\; ) \mid \mathbf{n}
  \end{array}
\]

e a entrada \haskell{n + n}. O analisador LL(1)
mantém uma \textbf{pilha} de símbolos que representa o que ainda
falta reconhecer. Inicialmente a pilha contém o símbolo inicial
seguido de \(\$\). A tabela a seguir mostra os passos para reconhecer
a string \haskell{n + n}:

\begin{table}[ht]
  \centering
  \begin{tabular}{lll}
    \toprule
    Pilha & Entrada & Ação \\
    \midrule
    $E\;\$$ & $\mathbf{n}\;+\;\mathbf{n}\;\$$ & Consultar $M[E,
    \mathbf{n}]$: aplicar $E \to T\;E'$ \\
    $T\;E'\;\$$ & $\mathbf{n}\;+\;\mathbf{n}\;\$$ & Consultar $M[T,
    \mathbf{n}]$: aplicar $T \to F\;T'$ \\
    $F\;T'\;E'\;\$$ & $\mathbf{n}\;+\;\mathbf{n}\;\$$ & Consultar
    $M[F, \mathbf{n}]$: aplicar $F \to \mathbf{n}$ \\
    $\mathbf{n}\;T'\;E'\;\$$ & $\mathbf{n}\;+\;\mathbf{n}\;\$$ &
    Casar: consumir $\mathbf{n}$ \\
    $T'\;E'\;\$$ & $+\;\mathbf{n}\;\$$ & Consultar $M[T', +]$:
    aplicar $T' \to \lambda$ \\
    $E'\;\$$ & $+\;\mathbf{n}\;\$$ & Consultar $M[E', +]$: aplicar
    $E' \to +\;T\;E'$ \\
    $+\;T\;E'\;\$$ & $+\;\mathbf{n}\;\$$ & Casar: consumir $+$ \\
    $T\;E'\;\$$ & $\mathbf{n}\;\$$ & Consultar $M[T, \mathbf{n}]$:
    aplicar $T \to F\;T'$ \\
    $F\;T'\;E'\;\$$ & $\mathbf{n}\;\$$ & Consultar $M[F,
    \mathbf{n}]$: aplicar $F \to \mathbf{n}$ \\
    $\mathbf{n}\;T'\;E'\;\$$ & $\mathbf{n}\;\$$ & Casar: consumir
    $\mathbf{n}$ \\
    $T'\;E'\;\$$ & $\$$ & Consultar $M[T', \$]$: aplicar $T' \to \lambda$ \\
    $E'\;\$$ & $\$$ & Consultar $M[E', \$]$: aplicar $E' \to \lambda$ \\
    $\$$ & $\$$ & \textbf{Aceitar} \\
    \bottomrule
  \end{tabular}
  \caption{Simulação do analisador preditivo para $\mathbf{n}+\mathbf{n}$}
\end{table}

Em cada passo, o analisador observa o topo da pilha e o símbolo
corrente da entrada:

\begin{itemize}
  \item
    Se o topo é um \textbf{terminal} igual ao símbolo corrente, ele é
    consumido da entrada e removido da pilha (casamento).
  \item
    Se o topo é um \textbf{não-terminal} \(A\), consulta-se \(M[A,
    a]\) para obter a produção a aplicar. O não-terminal é removido e
    os símbolos do corpo da produção são empilhados em ordem (o
    primeiro símbolo da produção fica no topo).
  \item
    Se o topo e a entrada são ambos \(\$\), a entrada foi aceita.
  \item
    Qualquer outra situação é um erro sintático.
\end{itemize}

Note que o problema central em um analisador LL(1) consiste em
construir a tabela. Esse será o assunto das próximas seções.

\section{Construção da tabela LL(1)}

Para construir a tabela de parsing, precisamos de três conceitos,
definidos a seguir.

\subsection{Anulável}

Um não-terminal \(A\) é \textbf{anulável} se pode derivar a cadeia
vazia em um número qualquer de passos:

\[
  A \Rightarrow^* \lambda
\]

Uma sequência de símbolos \(\alpha = X_1\, X_2\, \cdots\, X_n\) é
anulável se \textbf{todos} os \(X_i\) são anuláveis. A sequência
vazia é trivialmente anulável.

\subsection{Conjunto FIRST}

O conjunto \(\mathrm{FIRST}(\alpha)\) de uma sequência de símbolos
\(\alpha\) é o conjunto de terminais que podem iniciar uma cadeia
derivada de \(\alpha\):

\[
  \mathrm{FIRST}(\alpha) = \{\, a \in \Sigma \mid \alpha
  \Rightarrow^* a\,\beta \,\}
\]

Adicionalmente, incluímos \(\lambda\) em \(\mathrm{FIRST}(\alpha)\)
se \(\alpha\) é anulável.

As regras de cálculo para uma sequência \(X_1\, X_2\, \cdots\, X_n\) são:

\begin{enumerate}
  \item
    Se a sequência é vazia: \(\mathrm{FIRST}(\varepsilon) = \{\lambda\}\).
  \item
    Se \(X_1\) é um terminal \(a\): \(\mathrm{FIRST}(X_1 \cdots X_n) = \{a\}\).
  \item
    Se \(X_1\) é um não-terminal:

    \begin{itemize}
      \item
        Incluir \(\mathrm{FIRST}(X_1) - \{\lambda\}\).
      \item
        Se \(\lambda \in \mathrm{FIRST}(X_1)\), incluir também
        \(\mathrm{FIRST}(X_2 \cdots X_n)\).
    \end{itemize}
\end{enumerate}

Os conjuntos \(\mathrm{FIRST}\) para cada não-terminal são calculados
iterativamente até atingir um ponto fixo, isto é, até não haver
modificações no conjunto.

\subsection{Conjunto FOLLOW}

O conjunto \(\mathrm{FOLLOW}(A)\) de um não-terminal \(A\) é o
conjunto de terminais (e o marcador \(\$\)) que podem aparecer
imediatamente à direita de \(A\) em alguma forma sentencial:

\[
  \mathrm{FOLLOW}(A) = \{\, a \in \Sigma \cup \{\$\} \mid S
  \Rightarrow^* \alpha\, A\, a\, \beta \,\}
\]

As regras de cálculo são:

\begin{enumerate}
  \item
    Se \(A\) é o símbolo inicial: \(\$ \in \mathrm{FOLLOW}(A)\).
  \item
    Para cada regra \(B \to \alpha\, A\, \beta\):
    \begin{itemize}
      \item
        Adicionar \(\mathrm{FIRST}(\beta) - \{\lambda\}\) a
        \(\mathrm{FOLLOW}(A)\).
      \item
        Se \(\lambda \in \mathrm{FIRST}(\beta)\), adicionar
        \(\mathrm{FOLLOW}(B)\) a \(\mathrm{FOLLOW}(A)\).
    \end{itemize}
\end{enumerate}

Os conjuntos \(\mathrm{FOLLOW}\) também são calculados iterativamente
até ponto fixo.

\subsection{Algoritmo para construção da tabela de parsing}

A tabela \(M\) é uma matriz indexada por pares \((A, a)\), onde \(A\)
é um não-terminal e \(a\) é um terminal ou \(\$\). Cada entrada
contém uma produção ou indica erro. O algoritmo de construção é:
\begin{algorithm}
  \begin{algorithmic}[1]
    \ForAll{produção $A \to \alpha$ da gramática}
    \ForAll{terminal $a \in \mathrm{FIRST}(\alpha) - \{\lambda\}$}
    \State $M[A, a] \gets A \to \alpha$
    \EndFor
    \If{$\lambda \in \mathrm{FIRST}(\alpha)$}
    \ForAll{$b \in \mathrm{FOLLOW}(A)$}
    \State $M[A, b] \gets A \to \alpha$
    \EndFor
    \EndIf
    \EndFor
    \State Todas as entradas não preenchidas indicam \textbf{erro}
  \end{algorithmic}
  \caption{Construção da tabela de análise preditiva $M$}
\end{algorithm}

A gramática é \textbf{LL(1)} se e somente se nenhuma entrada da
tabela receber mais de uma produção. Quando isso ocorre, diz-se que
há um \textbf{conflito} na tabela.

\section{Algoritmo de Análise Sintática LL(1)}

Com a tabela construída, o algoritmo de parsing é:

\begin{algorithm}
  \begin{algorithmic}[1]
    \Require tabela $M$, símbolo inicial $S$, sequência de tokens $w\$$
    \State pilha $\gets [S,\,\$]$
    \Loop
    \State Seja $X$ o topo da pilha e $a$ o símbolo corrente da entrada
    \If{$X = \$$ e $a = \$$}
    \State \textbf{aceitar} e parar
    \ElsIf{$X$ é um terminal}
    \If{$X = a$}
    \State Desempilhar $X$ e avançar na entrada
    \Else
    \State \textbf{erro}
    \EndIf
    \ElsIf{$X$ é um não-terminal}
    \If{$M[X, a] = X \to \alpha$}
    \State Substituir $X$ na pilha pelos símbolos de $\alpha$
    \Statex
    \hspace{\algorithmicindent}\hspace{\algorithmicindent}\hspace{\algorithmicindent}(empilhados
    da direita para a esquerda)
    \Else
    \State \textbf{erro}
    \EndIf
    \EndIf
    \EndLoop
  \end{algorithmic}
  \caption{Analisador sintático preditivo dirigido por tabela}
\end{algorithm}

A sequência de produções aplicadas corresponde à \textbf{derivação
mais à esquerda} da entrada.

\section{Exemplo de Construção de Tabela}

\subsection{Gramática LL(1): Expressões Aritméticas}

Retomamos a gramática de expressões apresentada no exemplo informal:

\[
  \begin{array}{lcl}
    E  & \to & T\; E' \\
    E' & \to & +\; T\; E' \mid \lambda \\
    T  & \to & F\; T' \\
    T' & \to & *\; F\; T' \mid \lambda \\
    F  & \to & (\; E\; ) \mid \mathbf{n}
  \end{array}
\]

\textbf{Conjuntos FIRST:}

\begin{table}[ht]
  \centering
  \begin{tabular}{ll}
    \toprule
    Não-terminal & FIRST \\
    \midrule
    $E$  & $\{(,\;\mathbf{n}\}$ \\
      $E'$ & $\{+,\;\lambda\}$ \\
      $T$  & $\{(,\;\mathbf{n}\}$ \\
        $T'$ & $\{*,\;\lambda\}$ \\
        $F$  & $\{(,\;\mathbf{n}\}$ \\
          \bottomrule
        \end{tabular}
        \caption{Conjuntos FIRST}
      \end{table}

      \textbf{Conjuntos FOLLOW:}

      \begin{table}[ht]
        \centering
        \begin{tabular}{ll}
          \toprule
          Não-terminal & FOLLOW \\
          \midrule
        $E$  & $\{\$,\;)\}$ \\
      $E'$ & $\{\$,\;)\}$ \\
    $T$  & $\{+,\;\$,\;)\}$ \\
  $T'$ & $\{+,\;\$,\;)\}$ \\
$F$  & $\{*,\;+,\;\$,\;)\}$ \\
\bottomrule
\end{tabular}
\caption{Conjuntos FOLLOW}
\end{table}

\textbf{Tabela de parsing LL(1):}

\begin{table}[ht]
\centering
\begin{tabular}{lllllll}
\toprule
& $\mathbf{n}$ & $+$ & $*$ & $($ & $)$ & $\$$ \\
\midrule
$E$  & $T\;E'$      & $\varnothing$& $\varnothing$         & $T\;E'$      & $\varnothing$
& $\varnothing$       \\
$E'$ & $\varnothing$          & $+\;T\;E'$  & $\varnothing$         & $\varnothing$          &
$\lambda$  & $\lambda$  \\
$T$  & $F\;T'$      & $\varnothing$         & $\varnothing$         & $F\;T'$      &
$\varnothing$        & $\varnothing$       \\
$T'$ & $\varnothing$          & $\lambda$    & $*\;F\;T'$  & $\varnothing$          &
$\lambda$  & $\lambda$  \\
$F$  & $\mathbf{n}$ & $\varnothing$         & $\varnothing$         & $(\;E\;)$    &
$\varnothing$        & $\varnothing$       \\
\bottomrule
\end{tabular}
\caption{Tabela de análise preditiva $M$}
\end{table}

Cada entrada tem \textbf{no máximo uma produção}, confirmando que a
gramática é LL(1). A seguir, apresentamos alguns exemplos de
gramáticas que não são LL(1).

\subsection{Gramática não-LL(1): Conflito FIRST/FIRST}

Considere a gramática para comandos condicionais com o problema do
\emph{dangling else}:

\[
\begin{array}{lcl}
S & \to & \mathbf{if}\; E\; \mathbf{then}\; S\; \mathbf{else}\; S \\
& \mid & \mathbf{if}\; E\; \mathbf{then}\; S \\
& \mid & \mathbf{other}
\end{array}
\]

Para o não-terminal \(S\) com lookahead \(\mathbf{if}\):

\begin{itemize}
\item
\(\mathrm{FIRST}(\mathbf{if}\;E\;\mathbf{then}\;S\;\mathbf{else}\;S)
= \{\mathbf{if}\}\)
\item
\(\mathrm{FIRST}(\mathbf{if}\;E\;\mathbf{then}\;S) = \{\mathbf{if}\}\)
\end{itemize}

Ambas as produções contribuem para \(M[S, \mathbf{if}]\), gerando um
\textbf{conflito FIRST/FIRST}. A gramática não é LL(1) porque, com
apenas um símbolo de lookahead, não é possível distinguir dentre as
duas alternativas.

\subsection{Gramática não-LL(1): Conflito FIRST/FOLLOW}

Outro tipo comum de conflito ocorre quando uma produção pode derivar
\(\lambda\) e os conjuntos FIRST e FOLLOW do não-terminal não são
disjuntos. Considere:

\[
\begin{array}{lcl}
S & \to & A\; \mathbf{a} \\
A & \to & \mathbf{a} \mid \lambda
\end{array}
\]

Aqui \(\mathrm{FIRST}(A) = \{\mathbf{a}, \lambda\}\) e
\(\mathrm{FOLLOW}(A) = \{\mathbf{a}\}\). Para a entrada \(M[A,
\mathbf{a}]\) devemos registrar tanto \(A \to \mathbf{a}\) (pois
\(\mathbf{a} \in \mathrm{FIRST}(A)\)) quanto \(A \to \lambda\) (pois
\(\mathbf{a} \in \mathrm{FOLLOW}(A)\)). Há um \textbf{conflito
FIRST/FOLLOW} e, portanto, a gramática não é LL(1).

A seguir, vamos discutir a implementação deste algoritmo em Haskell.

\section{A Implementação em Haskell}

A implementação está dividida em três módulos principais:
\begin{flushleft}
\verb|Parsing.Grammar|,\\
\verb|Parsing.LL.BuildTable| e\\
\verb|Parsing.LL.Parser|.
\end{flushleft}

\subsection{Representação da Gramática}

Uma gramática é representada como um mapa de não-terminais para
listas de produções:

\begin{minted}{haskell}
data Symbol = Terminal String | NonTerminal String

type Production = [Symbol]

type Grammar = Map String [Production]
\end{minted}

O terminal especial $\lambda$ representa a
produção vazia. Funções auxiliares como
\haskell{nonTerminals},
\haskell{terminals},
\haskell{isLambda} e
\haskell{symbolString} facilitam a manipulação
das estruturas.

\subsection{Conjuntos FIRST (Parsing.First)}

Os conjuntos FIRST são computados por iteração até ponto fixo, usando
a função \haskell{fixedPoint} do módulo
\verb|Utils.Fixpoint|:

\begin{minted}{haskell}
type FirstSet = Map String [String]

computeFirstSets :: Grammar -> FirstSet
computeFirstSets g = fixedPoint step initial
  where
    initial = Map.fromList [(nt, []) | nt <- nonTerminals g]
    step cur = Map.mapWithKey update cur
      where
        update nt _ = unions $ map (firstForSequence g cur)
                                   (fromMaybe [] (Map.lookup nt g))
\end{minted}

A função \haskell{firstForSequence} implementa
diretamente as regras da definição:

\begin{minted}{haskell}
firstForSequence :: Grammar -> FirstSet -> [Symbol] -> [String]
firstForSequence _ _ [] = ["λ"]
firstForSequence g fs (sym:syms)
    | isTerminal sym =
        if isLambda sym
        then firstForSequence g fs syms
        else [symbolString sym]
    | otherwise =
        let nt         = symbolString sym
            firstOfSym = fromMaybe [] (Map.lookup nt fs)
            withoutLam = firstOfSym \\ ["λ"]
        in if "λ" `elem` firstOfSym
           then withoutLam `union` firstForSequence g fs syms
           else withoutLam
\end{minted}

O caso base retorna \haskell{["λ"]} para a sequência
vazia. Para um terminal, retorna o próprio terminal (ou continua com
o resto da sequência se for \haskell{λ}. Para um
não-terminal, inclui seu FIRST sem \haskell{λ} e, se
ele é anulável, continua com o restante da sequência.

\subsection{Conjuntos FOLLOW (Parsing.Follow)}

Os conjuntos FOLLOW também são calculados por ponto fixo. O símbolo
inicial começa com \haskell{"$"} em seu conjunto
FOLLOW, e os demais começam vazios:

\begin{minted}{haskell}
computeFollowSets :: Grammar -> FirstSet -> FollowSet
computeFollowSets g fs = fixedPoint step initial
  where
    startSymbol = head (nonTerminals g)
    initial = Map.fromList
        [ (nt, if nt == startSymbol then ["$"] else [])
        | nt <- nonTerminals g ]
    step cur = foldl update cur allProductions
\end{minted}

Para cada ocorrência de um não-terminal \(B\) numa produção \(A \to
\alpha\,B\,\beta\), o código calcula \(\mathrm{FIRST}(\beta)\) e
adiciona à \(\mathrm{FOLLOW}(B)\); se \(\beta\) é anulável, adiciona
também \(\mathrm{FOLLOW}(A)\):

\begin{minted}{haskell}
updatePos lhs' follow' pos
  | NonTerminal b <- prod !! pos =
      let beta      = drop (pos + 1) prod
          firstBeta = firstForSequence g fs beta
          followB   = fromMaybe [] (Map.lookup b follow')
          newFollow  = followB
                       `union` (firstBeta \\ ["λ"])
                       `union` (if "λ" `elem` firstBeta
                                then fromMaybe [] (Map.lookup lhs' cur)
                                else [])
      in Map.insert b newFollow follow'
  | otherwise = follow'
\end{minted}

\subsection{Construção da Tabela (Parsing.LL.BuildTable)}

A tabela é um mapa de pares \haskell{(String,String)} para
\haskell{ParseAction}:

\begin{minted}{haskell}
type ParseTable = Map (String, String) ParseAction

data ParseAction
  = Derive Production
  | Accept
  | Error
\end{minted}

A tabela é inicializada com \haskell{Error} em todas
as entradas \haskell{(nt, t)} para cada par
(não-terminal, terminal). Em seguida, cada produção preenche as
entradas correspondentes:

\begin{minted}{haskell}
buildParseTable :: Grammar -> ParseTable
buildParseTable g =
    foldl (addProduction g firstSets followSets) (initialTable g) (allProductions g)
  where
    firstSets  = computeFirstSets g
    followSets = computeFollowSets g firstSets
\end{minted}

A função \haskell{addProduction} determina os
terminais para os quais a produção deve ser registrada
e detecta conflitos:

\begin{minted}{haskell}
addProduction :: Grammar -> FirstSet -> FollowSet
              -> ParseTable -> (String, Production) -> ParseTable
addProduction g firstSets followSets table (lhs, prod) =
    foldl insertEntry table entries
  where
    firstAlpha = firstForSequence g firstSets prod
    followLhs  = fromMaybe [] (Map.lookup lhs followSets)

    entries :: [String]
    entries = (firstAlpha \\ ["λ"])
           ++ if "λ" `elem` firstAlpha then followLhs else []

    insertEntry tbl a =
      case Map.lookup (lhs, a) tbl of
        Just Error -> Map.insert (lhs, a) (Derive prod) tbl
        Just _     -> error $ "Grammar is not LL(1): conflict at ("
                           ++ lhs ++ ", " ++ a ++ ")"
        Nothing    -> tbl
\end{minted}

A lista \haskell{entries} reúne exatamente os
terminais do algoritmo de construção: \(\mathrm{FIRST}(\alpha)
\setminus \{\lambda\}\) mais, se \(\alpha\) é anulável,
\(\mathrm{FOLLOW}(A)\). Se uma entrada já tem uma ação diferente de
\haskell{Error}, a gramática não é LL(1) e uma
exceção é lançada imediatamente.

\subsection{O Parser (Parsing.LL.Parser)}

O analisador é uma função pura que recebe a gramática, a tabela e a
sequência de tokens, e retorna as produções aplicadas ou uma
mensagem de erro:

\begin{minted}{haskell}
data ParseResult
  = ParseOk [Production]
  | ParseError String

parse :: Grammar -> ParseTable -> [String] -> ParseResult
parse grammar table tokens =
    go initialStack (tokens ++ ["$"]) []
  where
    startSym     = head (nonTerminals grammar)
    initialStack = [NonTerminal startSym, Terminal "$"]
\end{minted}

O loop principal tem quatro casos, correspondendo diretamente ao algoritmo:

\begin{minted}{haskell}
    go (Terminal "$" : _) ("$" : _) acc =
        ParseOk (reverse acc)

    go (Terminal t : stack') (a : inp') acc
        | t == a    = go stack' inp' acc
        | otherwise = ParseError $
            "Syntax error: expected '" ++ t ++ "', found '" ++ a ++ "'"

    go (NonTerminal nt : stack') inp@(a : _) acc =
        case Map.lookup (nt, a) table of
            Just (Derive prod) ->
                let pushed = filter (not . isLambda) prod
                in go (pushed ++ stack') inp (prod : acc)
            Just Accept ->
                ParseOk (reverse acc)
            _ ->
                ParseError $
                    "Syntax error: no rule for (" ++ nt ++ ", " ++ a ++ ")"

    go stack' inp' _ =
        ParseError "Syntax error: unexpected state"
\end{minted}

Quando uma produção \(\lambda\) é aplicada, os símbolos
\haskell{λ} são filtrados antes de ser empilhados
(\haskell{filter (not . isLambda)}), pois \(\lambda\)
representa a cadeia vazia e não deve ocupar espaço na pilha. O
acumulador \haskell{acc} coleta as produções em ordem
reversa; ao aceitar, \haskell{reverse acc} devolve a
derivação mais à esquerda na ordem correta.

\section{Conclusão}

Neste capítulo, estudamos a análise sintática LL(1) do ponto de vista
teórico e prático.

Do lado teórico, definimos os conceitos de \textbf{anulável},
\textbf{FIRST} e \textbf{FOLLOW} e vimos como eles se combinam para
construir a tabela de parsing preditivo. A tabela encapsula toda a
decisão do analisador: dado um não-terminal e um símbolo de
lookahead, ela indica exatamente qual produção aplicar, ou que a
entrada é inválida.

Do lado prático, a implementação em Haskell mostrou que o algoritmo
se traduz diretamente em código. Os conjuntos FIRST e FOLLOW são
calculados por iteração até ponto fixo, a tabela é preenchida com um
único percurso sobre as produções da gramática, e o parser em si é
uma função recursiva com três casos simples.

Uma limitação importante da abordagem LL(1) é que nem toda gramática
livre de contexto pode ser analisada desta forma: gramáticas com
\textbf{recursão à esquerda} ou que não sejam \textbf{fatoradas à
esquerda} geram conflitos na tabela. Nesses casos, é necessário
transformar a gramática ou usar um algoritmo mais poderoso, como
os analisadores \textbf{LR} que estudaremos nos próximos capítulos.
